\DocumentMetadata{
  pdfversion=2.0,
  pdfstandard=ua-2,
  lang=en-US,
  tagging=on,
  tagging-setup={math/setup=mathml-SE,tabsorder=structure}
}
\documentclass[11pt,letterpaper]{article}
\usepackage[margin=0.68in,includefoot]{geometry}
\usepackage{newcomputermodern}
\usepackage{amsmath,amscd,mathtools}
\providecommand{\square}{\mdlgwhtsquare}
\usepackage[hidelinks]{hyperref}
\hypersetup{
  pdftitle={Differential Geometry PhD Exam, May 1992},
  pdfauthor={University of Florida Department of Mathematics},
  pdfsubject={Graduate mathematics examination},
  pdfdisplaydoctitle=true,
  bookmarksopen=true
}
\setcounter{secnumdepth}{0}
\setlength{\parindent}{0pt}
\setlength{\parskip}{0.28em}
\setlength{\emergencystretch}{3em}
\setlength{\leftmargini}{2.15em}
\setlength{\leftmarginii}{2.65em}
\setlength{\leftmarginiii}{3em}
\renewcommand{\labelenumi}{\textbf{\theenumi.}}
\pagestyle{empty}
\raggedbottom
\allowdisplaybreaks
\newenvironment{examproblems}[1]{%
  \begin{enumerate}%
  \setcounter{enumi}{#1}%
  \setlength{\topsep}{0.35em}%
  \setlength{\partopsep}{0pt}%
  \setlength{\itemsep}{0.48em}%
  \setlength{\parsep}{0.18em}%
}{\end{enumerate}}
\newenvironment{examsubparts}{%
  \begin{enumerate}%
  \setlength{\topsep}{0.18em}%
  \setlength{\partopsep}{0pt}%
  \setlength{\itemsep}{0.16em}%
  \setlength{\parsep}{0.1em}%
}{\end{enumerate}}
\newenvironment{examinstructions}{%
  \par\begingroup\itshape
}{%
  \par\endgroup\vspace{0.18em}
}
\makeatletter
\renewcommand\section{\@startsection{section}{1}{0pt}%
  {-1.25ex plus -.25ex minus -.1ex}%
  {0.55ex plus .1ex}%
  {\normalfont\large\bfseries}}
\renewcommand{\@maketitle}{%
  \newpage\null\vskip -1.2em%
  {\footnotesize\bfseries
    \href{https://www.ufl.edu/}{University of Florida}, \href{https://math.ufl.edu/}{Department of Mathematics}\par}%
  \vskip 0.18em%
  \tagstructbegin{tag=Title}%
    {\LARGE\bfseries\href{https://gma.math.ufl.edu/exams/differential-geometry-phd/}{Differential Geometry PhD Exam}, May 1992\par}%
  \tagstructend%
  \vskip 0.35em%
  \tagmcbegin{artifact}%
    \hrule height 1pt%
  \tagmcend%
  \vskip 0.55em%
}
\makeatother
\title{Differential Geometry PhD Exam, May 1992}
\author{}
\date{}
\begin{document}
\maketitle
\thispagestyle{empty}
\begin{examinstructions}
Answer three questions from Section I and each question in Section II.
\end{examinstructions}
\section{SECTION I}
\begin{examproblems}{0}
\item
State the preimage (or regular value) theorem. Let~\(M\) be the vector space of \(m\times m\) real matrices and~\(S\) its subspace of symmetric matrices. By considering the function 
\[
F:M\longrightarrow S:A\longmapsto A^tA
\]
 or otherwise, show that the orthogonal group \(O(m)\) is naturally a submanifold of~\(M\) having dimension \(\frac12m(m-1)\). With which subspace of~\(M\) is \(T_I O(m)\) naturally identified?
\item
What are the integral curves of a vector field. What does it mean to say that a vector field is complete? Show that \(\xi=y^2\frac{\partial}{\partial x}\) and \(\eta=x^2\frac{\partial}{\partial y}\) are complete vector fields on \(\mathbb R^2\) but that \(\xi+\eta\) is not complete. Can the Lie bracket of two different incomplete vector fields be complete? Justify.
\item
How is the exterior derivative \(d\omega\in\Omega^{k+1}(M)\) of \(\omega\in\Omega^k(M)\) defined? Let~\(\omega\) be a volume form on the orientable manifold~\(M\) and let~\(\xi\) be a vector field on~\(M\). Show that \(L_\xi\omega\) is exact and that \(L_\xi\omega=f_\xi\omega\) for some smooth function \(f_\xi\in C(M)\) on~\(M\). Find an explicit formula for \(f_\xi\) if 
\[
\xi=a_1\frac{\partial}{\partial x_1}+\cdots+a_m\frac{\partial}{\partial x_m}
\]
 on \(M=\mathbb R^m\) with \(a_1,\ldots,a_m\in C(\mathbb R^m)\) and \(\omega=dx_1\wedge\cdots\wedge dx_m\) in the usual coordinates. Hence suggest a suitable name for \(f_\xi\) in general.
\item
Let \(\omega=x\,dy-y\,dx\) be the standard one-form on the unit circle \(S^1\). Let \(f:\mathbb R\longrightarrow S^1\) be the inverse of stereographic projection from \((0,1)\). Compute \(f^*\omega\) and hence show that 
\[
\int_{S^1}\omega=2\pi.
\]
 How does this calculation establish that~\(\omega\) is not exact on \(S^1\)? Is \(f^*\omega\) exact on \(\mathbb R\)? Why?
\end{examproblems}
\section{SECTION II}
\begin{examproblems}{4}
\item
Explain what is meant by a Lie group~\(G\), its Lie algebra \(\mathfrak g\), and the adjoint representation \(\operatorname{Ad}\) of~\(G\) on \(\mathfrak g\). How does the adjoint representation of \(\operatorname{Gl}(m,\mathbb R)\) appear in terms of the natural identifications? Show that if~\(G\) is connected then the kernel of the adjoint representation coincides with the centre of~\(G\). It may be assumed that if \(t\in\mathbb R\), \(g\in G\), and \(\xi\in\mathfrak g\), then 
\[
\exp(t\operatorname{Ad}_g\xi)=g(\exp t\xi)g^{-1}.
\]
\item
Let \((M,\omega)\) be a symplectic manifold. How is the Hamiltonian vector field \(\xi_H\) of \(H\in C(M)\) defined? How is the Poisson bracket defined on \(C(M)\)? Show that if \((p_1,\ldots,p_m,q_1,\ldots,q_m)\) are (local) symplectic coordinates on~\(M\), then the integral curves of \(\xi_H\) satisfy the Hamilton equations 
\[
\dot q_j=\frac{\partial H}{\partial p_j},\qquad \dot p_j=-\frac{\partial H}{\partial q_j}.
\]
 Show also that the Poisson centralizer of \(H\in C(M)\) is precisely the set of all functions on~\(M\) constant along integral curves of \(\xi_H\), and suggest a mechanical interpretation of this result.
\item
Define a geodesic in terms of the Levi-Civita connection~\(\nabla\) on the Riemannian manifold \((M,g)\) and give the form taken by the geodesic equation in local coordinates. For the metric~\(g\) on the open upper half-plane~\(M\) having \(\left\{y\frac{\partial}{\partial x},y\frac{\partial}{\partial y}\right\}\) as orthonormal vector fields, show that the nonzero Christoffel symbols are 
\[
\Gamma^y_{yy}=\Gamma^x_{xy}=\Gamma^x_{yx}=-\Gamma^y_{xx}=-\frac1y.
\]
 Hence show that vertical upper half-lines and upper semicircles centered on the~\(x\)-axis are geodesics. It may be assumed that if \(\xi,\eta,\zeta\) are vector fields on~\(M\), then 
\[
\begin{aligned}
2g(\nabla_\xi\eta,\zeta)={}&\xi\cdot g(\eta,\zeta)+\eta\cdot g(\xi,\zeta)-\zeta\cdot g(\xi,\eta)\\
&-g(\xi,[\eta,\zeta])-g(\eta,[\xi,\zeta])+g(\zeta,[\xi,\eta]).
\end{aligned}
\]
\end{examproblems}
\end{document}
